\subsection{Cross-Participant 200-ms Prediction}
\label{subsec:unseen-200ms}

We separately examined cross-participant generalization at 200 ms using
50 open-hand and 50 closed-hand trials from an unseen participant, without
calibration or parameter updates. This horizon was selected because it is the
end of the densely supervised trajectory window and is therefore predicted by
the same type of head under the same target definition for every model. It
also provides a short anticipatory interval that can be consumed by a
downstream controller while remaining sufficiently close to the observed
motion for reliable evaluation. Restricting this comparison to 200 ms avoids
confounding the result with the proposed model's separate long-horizon intent
decoder.

\begin{table}[t]
 \centering
 \caption{Unseen-participant Cartesian prediction error at 200 ms. The pooled
 result is weighted by the number of valid frames in each condition.}
 \label{tab:unseen-200ms}
 \small
 \setlength{\tabcolsep}{4pt}
 \begin{tabular}{lrrr}
  \hline
  Model & Open (cm) & Closed (cm) & Pooled (cm) \\
  \hline
  Proposed       & 7.70 & \textbf{8.37} & \textbf{8.04} \\
  IMU--GRU       & 7.01 & 9.57 & 8.33 \\
  sEMG+IMU--GRU  & 8.27 & 15.01 & 11.73 \\
  IMU--LSTM      & \textbf{6.97} & 10.49 & 8.78 \\
  sEMG+IMU--LSTM & 7.88 & 11.68 & 9.83 \\
  \hline
 \end{tabular}
\end{table}

The proposed model achieved the lowest pooled error of 8.04 cm, compared with
8.33 cm for the strongest conventional pooled baseline, IMU--GRU. Although the
IMU-only recurrent models were more accurate in the open condition, their
errors increased substantially in the closed condition. The proposed model
instead maintained similar performance across hand states, with an
open-to-closed difference of 0.67 cm, compared with 2.56 cm for IMU--GRU and
3.52 cm for IMU--LSTM. Its frozen state decoder also achieved 95.86\%
macro-F1 after the open and closed trials were combined. These observations
support the intended role of state conditioning: it reduces sensitivity to
changes in muscle recruitment and arm dynamics associated with the gripper
state while preserving an IMU-derived mechanical base.

Direct early fusion did not provide the same robustness. Its pooled 200-ms
errors increased to 11.73 cm for the GRU and 9.83 cm for the LSTM, indicating
that unconstrained sEMG features can introduce negative transfer under
participant shift. In contrast, the proposed bounded sEMG correction improved
the pooled result without allowing the muscle stream to replace the inertial
representation. The comparison therefore isolates the benefit of the proposed
fusion mechanism rather than merely demonstrating that an additional modality
was available. Because this experiment contains one unseen participant, the
differences are reported descriptively and require confirmation on a larger
cross-participant cohort.
